\subsubsection*{Part I: Sun-Synchronous Orbits}

\begin{description}
\item[(a)] Since the orbit is Sun-Synchronous, this means that the satellite's
  line of nodes must drift $360^\circ$ along one sidereal year. This ensures
  that the satellite will always pass over a given meridian at the same local
  mean solar time. Hence, \hfill(1.0 points)
  \[
    \dot\Omega = \frac{360^\circ}{T_s}
      = \frac{2\pi}{365.2564\cdot24\cdot60\cdot60}
      = 1.991\cdot10^{-7}\ \mathrm{rad/s}
  \]
  \hfill(2.0 points)

\item[(b)] After one orbit, the difference in longitude perceived by the
  satellite on the ground track is associated with the combined effects of the
  rotation of the Earth and nodal precession. This is expressed by
  \[
    \lambda_2 = \lambda_1 - \omega_\oplus\cdot T + \dot\Omega\cdot T
  \]
  \hfill(2.0 points)

  From the Ground Track, $\lambda_1 = 90^\circ$ and
  $\lambda_2 \approx -135^\circ$ for five orbits. Therefore, knowing the value
  of $\omega_\oplus$, the angular velocity of the rotation of our planet,
  \[
    T = \frac{1}{5}\cdot\frac{\lambda_1 - \lambda_2}{\omega_\oplus - \dot\Omega}
      = \frac{1}{5}\cdot
        \frac{225^\circ\cdot\pi/180^\circ}
             {7.292\cdot10^{-5} - 1.991\cdot10^{-7}}
      \approx 10800\ \mathrm{s} = 180\ \mathrm{min}
  \]
  \hfill(2.0 points)

  Although $\dot\Omega$ is significantly smaller than $\omega_\oplus$, it is
  conceptually important to include the nodal precession rate in the formula.

  The inclination of the orbit can be obtained directly from the Ground
  Track. Taking any of the two points of the orbit with the largest (absolute)
  value of latitude, one finds
  \[
    |\phi_{max}| \approx 55^\circ
  \]
  \hfill(2.0 points)

  Since $i > 90^\circ$, $|\phi_{max}|$ is the supplementary angle of the
  inclination:
  \[
    i = 180^\circ - 55^\circ = 125^\circ
  \]
  \hfill(2.0 points)

\item[(c)] Having calculated the period associated with the orbit, the
  semi-major axis is directly obtained from
  \[
    \frac{T^2}{a^3} = \frac{4\pi^2}{GM_\oplus}
    \implies
    a = \left(\frac{T^2 GM_\oplus}{4\pi^2}\right)^{1/3}
      = 1.06\cdot10^4\ \mathrm{km}
  \]
  \hfill(2.0 points)

\item[(d)] Since the orbit is sun-synchronous, the satellite completes an
  integer number of orbits during one solar day. Hence,
  \[
    N_S = \frac{24\cdot60}{180} = 8\ \text{orbits}
  \]
  \hfill(1.0 points)

\item[(e)] The satellite takes more than 2 orbits and less than 3 orbits to go
  from Poland to Brazil. Using as reference the point $P_o$ ahead of Macei\'o
  that is exactly 3 orbits after Chorz\'ow,
  \[
    \lambda_o = \lambda_C - 3\cdot45^\circ
  \]
  \[
    \phi_0 = \phi_C
  \]
  Using sine law in the following spherical triangle,

% FIGURE NEEDED: 3D sphere diagram (2024 theory solutions, page 33 of 40):
% Earth sphere with equatorial plane, the orbit great circle crossing the
% equator at the ascending node Omega, inclination marked as 180 deg - i,
% arcs theta_0 and theta_M along the orbit to points at latitudes phi_0 and
% phi_M, with x, y, z axes drawn.

  \[
    \sin\theta = \frac{\sin\phi}{\sin(180^\circ - i)}
  \]
  \hfill(2.0 points)
  \[
    \theta_0 = \arcsin\left(\frac{\sin\phi_o}{\sin i}\right) = 69.9^\circ
  \]
  \[
    \theta_M = \arcsin\left(\frac{\sin\phi_M}{\sin i}\right) = -12.0^\circ
  \]
  The angle the satellite traveled will depend on how many complete orbits it
  took and the angle from the final whole orbit to Macei\'o:
  \[
    \Delta\theta = n\cdot360^\circ - (\theta_0 - \theta_M)
  \]
  \hfill(3.0 points)

  by the graph, $n = 3$, which results in $\Delta\theta = 998.1^\circ$.
  \hfill(1.0 points)

  Knowing its orbit is circular,
  \[
    \Delta t = T\cdot\frac{\Delta\theta}{360^\circ} = 8.32\ \mathrm{h}
  \]
  \hfill(1.0 points)

  The local time at Chorz\'ow over which the satellite passes by will be,
  considering $\lambda_M = 35.7^\circ W = -35.7^\circ$,
  \[
    T_C = T_M - \Delta t + \omega_\oplus^{-1}(\lambda_C - \lambda_M)
  \]
  \hfill(3.0 points)
  \[
    \boxed{T_C \approx 7\mathrm{h}20\mathrm{min}}
  \]
  \hfill(1.0 points)
\end{description}

\subsubsection*{Part II: Tundra orbits}

\begin{description}
\item[(f)] For the inclination, we need to take the maximum latitude reached
  by the satellite, which gives us $\boxed{i = 63^\circ}$.
  \hfill(1.0 points)

  For the period, we have to realize that in one orbit, the Earth's rotation
  does not cause a shift in the orbit, so it must be geosynchronous, that is
  $T = T_\oplus = \boxed{1436\ \mathrm{min}}$. \hfill(1.0 points)

  For the semi-major axis, we apply Kepler's Third Law:
  \[
    \frac{a^3}{T^2} = \frac{GM}{4\pi^2}
    \rightarrow \boxed{a = 42{,}164\ \mathrm{km}}
  \]
  \hfill(2.0 points)

\item[(g)] We will first find the expression for the area divided by the
  semi-latus rectum as a function of eccentricity so then we can apply the
  Kepler's Second Law to find the required time. \hfill(2.0 points)

  To do this, we will use the projection of the area of a circle when it is
  inclined by an angle $\arccos\left(\frac{b}{a}\right)$, as shown in the
  figure below:

% FIGURE NEEDED: auxiliary-circle diagram (2024 theory solutions, page 35 of
% 40): circle with centre C, angle theta at the centre, points Q, I, T, P
% marking the sector cut off beyond the semi-latus rectum, hatched circular
% segment (light and dark blue) between the chord and the arc near P, angle
% alpha at Q.

  Thus, we will first calculate the hatched area of the circle figure
  $A_{circle}$
  \[
    A_{circle} = A_{sector} - A_{triangle}
      = \frac{\theta}{2\pi}\cdot\pi a^2 - \frac{ae\cdot b}{2}
  \]
  \hfill(2.0 points)

  We can also use that $\theta = \frac{\pi}{2} - \alpha$ and
  $\alpha = \sin^{-1}(e)$. Therefore:
  \[
    A_{circle} = \left(\frac{\pi\cdot a^2}{4}
      - \frac{\sin^{-1}(e)\cdot a^2}{2} - \frac{ab\cdot e}{2}\right)
  \]
  Now, we project this area,
  \[
    A_{ellipse} = A_{circle}\cdot\frac{b}{a}
  \]
  \[
    A_{ellipse} = \left(\frac{\pi ab}{4}
      - \frac{\sin^{-1}(e)\cdot ab}{2} - \frac{b^2 e}{2}\right)
  \]
  \hfill(4.0 points)

  Thus, the area covered by the satellite's position vector from one
  semi-latus rectum to the other will be given by twice the area found above,
  as we can see from the image, so that the area on the perigee side is given
  by
  \[
    A_{per} = \frac{\pi ab}{2} - b(a\sin^{-1}(e) + eb)
  \]
  For the apogee side, it is enough to know that
  $A_{ap} = \pi ab - A_{per}$, where $\pi ab$ is the area of the whole
  ellipse. Therefore, we have:
  \[
    A_{ap} = \frac{\pi ab}{2} + b\left(a\sin^{-1}(e) + eb\right)
  \]
  \hfill(2.0 points)

  Using Kepler's Second Law, we can write:
  \[
    \frac{A_{ap}}{A_{ellipse}} = \frac{T'}{T}
  \]
  \hfill(2.0 points)

  Using that $A_{ellipse} = \pi ab$ and $b = a\sqrt{1 - e^2}$, we have:
  \[
    \boxed{\frac{T'}{T} = \frac{1}{2} + \frac{\sin^{-1}(e)}{\pi}
      + \frac{e}{\pi}\cdot\sqrt{1 - e^2}}
  \]

  \textbf{Alternative Solution}

% FIGURE NEEDED: eccentric-anomaly diagram (2024 theory solutions, page 36 of
% 40): ellipse (trajectory) with auxiliary circle, centre O, foci F and F',
% point P on the ellipse below P' on the circle, radius a, focal distance ea,
% eccentric anomaly E at O and true anomaly theta at F, apex A. Caption cites
% https://en.wikipedia.org/wiki/Eccentric_anomaly as source.

  The relation between the true anomaly and the mean anomaly is given by
  \[
    M = E - e\cdot\sin(E)
  \]
  From the figure (improve figure later), we observe that the satellite has % sic: "(improve figure later)" left in the official solutions
  the following properties:
  \begin{itemize}
    \item At perigee:
      \begin{itemize}
        \item $E = 0$
        \item $\theta = 0$
      \end{itemize}
    \item At the semi-latus rectum:
      \begin{itemize}
        \item $E = \cos^{-1}(e)$
        \item $\theta = 90^\circ$
      \end{itemize}
  \end{itemize}
  \hfill(2.0 points)

  Therefore, at the moment of passage through the semi-latus rectum, knowing
  that $\sin(E) = \sqrt{1 - e^2}$, we have
  \[
    M = \cos^{-1}(e) - e\cdot\sqrt{1 - e^2}
  \]
  \hfill(4.0 points)

  If we also use the perigee as the initial reference for the mean anomaly, we
  can associate the time of passage from the perigee to a certain point in the
  orbit by
  \[
    t = \frac{M}{\omega}
  \]
  where $\omega$ is the mean angular velocity of the orbit, given by
  $\omega = \frac{2\pi}{T}$, with $T$ being the orbital period. Thus, the time
  elapsed from the perigee to the semi-latus rectum is given by
  \[
    t = T\frac{\cos^{-1}(e) - e\cdot\sqrt{1 - e^2}}{2\pi}
  \]
  \hfill(2.0 points)

  Finally, the time $T'$ spent in the northern hemisphere can be found as
  follows:
  \[
    T' + 2t = T \rightarrow T' = T - 2t
  \]
  Thus, the desired ratio $T'/T$ will be:
  \[
    \frac{T'}{T} = 1 - \frac{2t}{T}
  \]
  \hfill(2.0 points)

  Substituting the expression for $t$, we get:
  \[
    \frac{T'}{T} = 1 - \frac{\cos^{-1}(e) + e\cdot\sqrt{1 - e^2}}{\pi}
  \]
  If we use the relation $\cos^{-1}(e) + \sin^{-1}(e) = \frac{\pi}{2}$, we
  finally obtain:
  \[
    \boxed{\frac{T'}{T} = \frac{1}{2} + \frac{\sin^{-1}(e)}{\pi}
      + \frac{e}{\pi}\cdot\sqrt{1 - e^2}}
  \]

\item[(h)] Between the two consecutive passages through the semi-latus rectum,
  we can use the following equation:
  \[
    \Delta\lambda = \pi - \omega_\oplus\cdot T'
  \]
  \hfill(2.0 points)

  Where $\Delta\lambda$ is the variation in longitude of the satellite between
  its passages through the semi-latus rectum, as shown in the image below.

% FIGURE NEEDED: annotated Tundra ground track (2024 theory solutions, page 38
% of 40): the figure-8 ground track of Figure 3 with the "First passage" and
% "Second passage" through the semi-latus rectum marked and the longitude
% difference between them indicated by red dashed lines.

  Therefore, we can write using the expression of the last item
  \[
    \Delta\lambda + \omega_\oplus\cdot T\cdot
      \left(\frac{1}{2} + \frac{\sin^{-1}(e)}{\pi}
      + \frac{e}{\pi}\cdot\sqrt{1 - e^2}\right) = \pi
  \]
  As $\omega_\oplus = \frac{2\pi}{T}$ (the orbit is geosynchronous), we will
  get:
  \begin{align*}
    \Delta\lambda + 2\cdot\sin^{-1}(e) + 2e\cdot\sqrt{1 - e^2} &= 0 \\
    \sin^{-1}(e) + e\cdot\sqrt{1 - e^2} &= -\frac{\Delta\lambda}{2}
  \end{align*}
  \hfill(2.0 points)

  Considering that $\Delta\lambda \approx -67.5^\circ$ from the graph (it is
  important to note that this value should be negative by the figure), we find
  from the above formula, by iteration, $\boxed{e \approx 0.3}$.
  \hfill(2.0 points)

  If we use the approximations $\sin(e) \approx e$ and $e^2 \ll 1$ given in
  the statement, we find:
  \[
    2e = -\frac{\Delta\lambda}{2}
  \]
  \hfill(3.0 points)

  which gives us an approximate answer $\boxed{e \approx 0.295 \approx 0.3}$
  \hfill(1.0 points)

\item[(i)] At the start and end points of retrograde motion, we have
  $\omega_{RA} = \omega_\oplus$, where $\omega_{RA}$ is the right ascension
  angular velocity of the satellite. \hfill(2.0 points)

  To find this velocity, we use the following spherical triangle:

% FIGURE NEEDED: sphere diagram (2024 theory solutions, page 39 of 40):
% spherical triangle formed by the orbit great circle and the equator, with
% inclination i at the node, arc theta_2 along the orbit, declination delta,
% and right ascension alpha along the equator.

  where $\theta_2 = \theta - 90^\circ$, and it is known that:
  \[
    \cos(\theta_2) = \cos(\delta)\cos(\alpha) \quad\text{(Cosine Law)}
  \]
  \hfill(1.0 points)
  \[
    \frac{\sin(\theta_2)}{\sin(\alpha)} = \frac{\cos(\delta)}{\cos(i)}
    \quad\text{(Sine Law)}
  \]
  \hfill(1.0 points)

  We substitute $\cos(\delta)$ from the first expression into the second
  expression, obtaining:
  \begin{align*}
    \frac{\sin(\theta_2)}{\sin(\alpha)}
      &= \frac{\cos(\theta_2)}{\cos(\alpha)\cdot\cos(i)} \\
    \tan(\theta_2)\cdot\cos(i) &= \tan(\alpha)
  \end{align*}
  We then derive the above expression with respect to time --- it is possible
  to skip the derivative step by just decomposing the satellite's angular
  velocity vector.
  \[
    \frac{\cos(i)}{\cos^2(\theta_2)}\,\dot\theta
      = \frac{\dot\alpha}{\cos^2(\alpha)}
  \]
  \hfill(2.0 points)

  Using the Cosine Law expression to replace $\cos(\alpha)$, we finally have:
  \[
    \omega_{RA} = \dot\alpha = \frac{\cos(i)}{\cos^2(\delta)}\,\dot\theta
  \]
  \hfill(2.0 points)

  We can then calculate the value of $\dot\theta$ as a function of the
  distance of the satellite from the center of the Earth, using the
  conservation of angular momentum:
  \[
    \dot\theta = \frac{2\pi}{T}\,\frac{a^2\sqrt{1 - e^2}}{r^2}
  \]
  \hfill(2.0 points)

  We still have the polar expression for the distance $r$ as a function of the
  polar radius, given by: % sic: "as a function of the polar radius"
  \[
    r = \frac{a(1 - e^2)}{1 + e\cdot\cos(\theta)}
  \]
  \hfill(2.0 points)

  Substituting everything into the expression that gives the condition for the
  start or end of retrograde motion, we find:
  \begin{align*}
    \frac{2\pi}{T}\,
      \frac{(1 + e\cos\theta)^2}{(1 - e^2)^{\frac{3}{2}}}\,
      \frac{\cos(i)}{\cos^2(\delta)}
      &= \omega_{earth} = \frac{2\pi}{T} \\
    \frac{(1 + e\cos\theta)^2}{(1 - e^2)^{\frac{3}{2}}}\,
      \frac{\cos(i)}{\cos^2(\delta)} &= 1
  \end{align*}
  \hfill(2.0 points)

  From the graph, we find the latitude --- or declination --- of the start or
  end of retrograde motion, given by $\delta \approx -32^\circ$.
  \hfill(2.0 points)

  Thus, substituting all the variables into the above equation and solving for
  $\theta$, we find: $\boxed{\theta = 54^\circ}$ and
  $\boxed{\theta = 306^\circ}$ as solutions. \hfill(2.0 points)

\item[(j)] For the figure-8 shape not to form, we need the inversion of motion
  in the northern hemisphere to occur exactly at the apogee position.
  Therefore, we will use exactly the same expression as the previous item, but
  with $\theta = 180^\circ$ and eccentricity as the unknown variable.
  \hfill(2.0 + 1.0 points)

  For the value of $\delta$, as the satellite is at the position of maximum
  declination, $\delta = i$. We have:
  \[
    \frac{(1 - e)^2}{(1 - e^2)^{\frac{3}{2}}}\,\frac{1}{\cos(i)} = 1
  \]
  \hfill(2.0 points)

  Solving the above equation, we obtain: $\boxed{e \approx 0.4}$
  \hfill(1.0 points)
\end{description}
